% Originally: content/week-12.tex, \section{Solutions}

\section{Solutions}
\label{sec:stokes_solutions}

\begin{exercisesolution}[Solution of \cref{ex:12.3}]
\textbf{1.} The domain is $\{(x,y,z) \in \mathbb{R}^3 : (x,y) \neq (0,0)\}$, i.e.\
$\mathbb{R}^3$ with the $z$-axis removed. It is connected (path-connected: any two points can
be joined by a path avoiding the axis), but not simply connected: $\gamma_1$, for instance, which
winds once around the axis, cannot be contracted to a point without crossing it. (It is also
not convex: the segment from $(-1,0,0)$ to $(1,0,0)$ leaves the domain.)

\textbf{2.} Writing $F = (F^1,F^2,F^3)$ with $F^1 = 2(xz+y)/(x^2+y^2)$, $F^2 =
2(yz-x)/(x^2+y^2)$, $F^3 = \log(x^2+y^2)$, a direct computation of each pair of mixed
partials gives
\[\partial_y F^3 = \partial_z F^2 = \frac{2y}{x^2+y^2}, \qquad
  \partial_z F^1 = \partial_x F^3 = \frac{2x}{x^2+y^2}, \qquad
  \partial_x F^2 = \partial_y F^1 = \frac{2(x^2-y^2-2xyz)}{(x^2+y^2)^2},\]
so $\curl F = (\partial_yF^3-\partial_zF^2,\ \partial_zF^1-\partial_xF^3,\ \partial_xF^2-
\partial_yF^1) = (0,0,0)$: $F$ is curl-free. On $\gamma_1$, $x=\cos t$, $y=\sin t$, $z=0$,
$x^2+y^2=1$, so $F(\gamma_1(t)) = (2\sin t,\,-2\cos t,\,0)$ and $\gamma_1'(t) = (-\sin t,
\cos t,0)$, giving
\[\int_{\gamma_1} F\cdot d\gamma_1 = \int_0^{2\pi} \big(2\sin t\cdot(-\sin t) + (-2\cos t)
  \cdot\cos t\big)\,dt = \int_0^{2\pi} -2\,dt = -4\pi.\]
Since this is nonzero on a closed loop, $F$ is \textbf{not} conservative. It is curl-free, but
its domain is not simply connected (part 1), so curl-free does not force conservative here.

\textbf{3.} Let $S := \{(x,y,z) : \sqrt{x^2+y^2} = 1+z/2,\ 0\leq z\leq2\}$, the cone swept out
by linearly interpolating the radius from $\gamma_1$ ($z=0$, radius $1$) to $\gamma_2$ ($z=2$,
radius $2$), parametrized by $H(s,t) := \big((1+s)\cos t,(1+s)\sin t,2s\big)$ for $(s,t) \in
[0,1]\times[0,2\pi]$. Since the radius $1+s \geq 1 > 0$ throughout, $S$ never meets the
$z$-axis and lies entirely in the domain of $F$; its oriented boundary (outward normal) is
$\gamma_2$ minus $\gamma_1$. Since $\curl F = 0$ on $S$ by part 2, \cref{thm:stokes_classical}
gives
\[\int_{\gamma_2} F\cdot d\gamma_2 - \int_{\gamma_1} F\cdot d\gamma_1 = \int_{\partial S}
  F\cdot ds = \int_S \curl F\cdot d\vec n = 0,\]
i.e.\ $\int_{\gamma_1} F\cdot d\gamma_1 = \int_{\gamma_2} F\cdot d\gamma_2$.

\textbf{4.} No. The disk $D := \{(x,y,z) : (x-6)^2+y^2 \leq 9,\ z=3\}$, oriented upward, has
oriented boundary exactly $\gamma_3$, and $D$ lies entirely in the domain of $F$ (it does not
meet the $z$-axis, since its center $(6,0,3)$ is at distance $6$ from the axis and its radius
is only $3$). By \cref{thm:stokes_classical} with $\curl F=0$,
\[\int_{\gamma_3} F\cdot d\gamma_3 = \int_{\partial D} F\cdot ds = \int_D \curl F\cdot d\vec n
  = 0.\]
Since $\int_{\gamma_1} F\cdot d\gamma_1 = -4\pi \neq 0 = \int_{\gamma_3} F\cdot d\gamma_3$, the
two do \emph{not} agree, and Stokes' theorem gives no reason to expect they would: unlike the
pair $(\gamma_1,\gamma_2)$, there is no surface joining $\gamma_1$ and $\gamma_3$ that stays
inside the domain, since $\gamma_1$ winds around the excluded axis and $\gamma_3$ does not.
\end{exercisesolution}

% Verified: solved independently, all four parts match Sol12_Analysis2_eng.pdf exactly.
% Note: the official solution's part 2 labels vanishing curl as "divergence-free" where it
% visibly means curl-free (computes curl F, not div F) --- a wording slip, not repeated here.

\begin{exercisesolution}[Solution of \cref{ex:12.4}]
\textbf{(a) False.} The Newton potential's gradient $F(x) = \nabla(1/\lvert x\rvert) =
-x\lvert x\rvert^{-3}$ is divergence-free on $U$ (it is $V_{-3}$ up to sign, from
\cref{ex:11.5}, with $n=3$), yet $I_r = \int_{\partial B_r} F\cdot n\,dS =
-\vol_2(\partial B_r) \neq 0$.

\textbf{(b) True.} Applying \cref{thm:gauss_divergence} on the shell $B_{r_2}\setminus
B_{r_1} \Subset U$ for $0<r_1<r_2$, whose boundary is $\partial B_{r_2}$ (outward) together
with $\partial B_{r_1}$ with reversed orientation,
\[0 = \int_{B_{r_2}\setminus B_{r_1}} \divg F\,dx = I_{r_2} - I_{r_1},\]
so $I_r$ is independent of $r$.

\textbf{(c) False.} Take $F(x) = x = \tfrac12\nabla\lvert x\rvert^2$, which is curl-free
(being a gradient). Then $\divg F \equiv 3$, so by \cref{thm:gauss_divergence}, $I_r =
\int_{B_r} 3\,dx = 3\vol_3(B_r) = 4\pi r^3$, which depends on $r$.

\textbf{(d) True.} If $F = \curl G$ then $\divg F = \divg(\curl G) = 0$ (a standard vector
identity), so part (b) applies and $I_r$ is independent of $r$. On the other hand,
\[\lvert I_r\rvert \leq \int_{\partial B_r} \lvert\langle F,n\rangle\rvert\,dS \leq
  \frac{C}{r}\,\vol_2(\partial B_r) = \frac{C}{r}\cdot4\pi r^2 = 4\pi C r \longrightarrow 0
  \quad\text{as } r\to0^+.\]
Since $I_r$ is constant in $r$ and this constant is a limit of values tending to $0$, $I_r=0$
for all $r>0$.
\end{exercisesolution}

% Verified: solved independently, all four parts match Sol12_Analysis2_eng.pdf.

\begin{exercisesolution}[Solution of \cref{ex:12.5}]
Writing $F = (F^1,F^2)$ with $F^1 = \lambda xe^y$, $F^2 = (y+1+x^2)e^y$, the integrability
condition is $\partial_yF^1 = \partial_xF^2$:
\[\partial_yF^1 = \lambda xe^y, \qquad \partial_xF^2 = 2xe^y.\]
Since $\mathbb{R}^2$ is convex (so integrability is equivalent to conservativity here), $F$
is conservative if and only if $\lambda=2$.

For $\lambda=2$, seek $f$ with $\partial_xf = F^1 = 2xe^y$: integrating in $x$,
$f(x,y) = x^2e^y + g(y)$ for some function $g$. Then $\partial_yf = x^2e^y+g'(y)$ must equal
$F^2 = (y+1+x^2)e^y$, so $g'(y) = (y+1)e^y$; integrating by parts, $\int(y+1)e^y\,dy =
ye^y - e^y + e^y = ye^y$ (up to an additive constant). Hence
\[f(x,y) = \boxed{\ (x^2+y)e^y\ } \qquad (\text{unique up to an additive constant}),\]
which one checks directly satisfies $\nabla f = F$.
\end{exercisesolution}

% Verified: solved independently, lambda=2 and potential (x^2+y)e^y match Sol12 exactly.

\begin{exercisesolution}[Solution of \cref{ex:12.8}]
The four quantities $I_b,I_t,I_l,I_r$ are already defined using the same outward normal field
$n$ on each of the four edges of $\partial Q$, so the total outer flux is simply their sum,
with no relative sign corrections needed: $\boxed{I_b+I_t+I_r+I_l}$, answer \textbf{(b)}.
\end{exercisesolution}

% Verified: solved independently, matches Sol12 answer (b).
% Correction: Claude Opus 5 (High) --- this comment belongs to ex:12.8 above. Flash appended its
% own solution between the two, leaving the comment claiming a block it had never checked.

\begin{exercisesolution}[Proof of \cref{ex:curl_identity_greens_formula}]
% Generator: Gemini 3.6 Flash (High)
% Correction: Claude Opus 5 (High) --- \operatorname{rot} -> \curl, J\Phi -> \Jac\Phi, and the line
% integrals rewritten in this chapter's own notation (\int_{\partial\Omega} F\cdot\tau\,dL, as in
% thm:green_theorem) rather than the physics d\mathbf{r}. Part (b) also justified properly: Flash
% asserted "d\Phi_2 = d\Psi_2 on \partial\Omega", which is false for the full differentials --- only
% the tangential derivative agrees, and that is all the line integral sees.
\begin{enumerate}[label=\textbf{(\alph*)}]
  \item Write $V = (V_1, V_2)\transp = (\Phi_1 \partial_1 \Phi_2, \Phi_1 \partial_2 \Phi_2)\transp$.
The scalar curl in the plane is $\curl V = \partial_1 V_2 - \partial_2 V_1$, so the product rule
gives
\begin{align*}
\curl V &= \partial_1 (\Phi_1 \partial_2 \Phi_2) - \partial_2 (\Phi_1 \partial_1 \Phi_2) \\
&= (\partial_1 \Phi_1 \partial_2 \Phi_2 + \Phi_1 \partial_1 \partial_2 \Phi_2) - (\partial_2 \Phi_1 \partial_1 \Phi_2 + \Phi_1 \partial_2 \partial_1 \Phi_2).
\end{align*}
Since $\Phi \in C^2(\overline{\Omega}, \mathbb{R}^2)$, Schwarz's lemma
(\cref{thm:schwarz_mixed_partials}) gives $\partial_1 \partial_2 \Phi_2 = \partial_2 \partial_1
\Phi_2$, and the two second-derivative terms cancel. Consequently,
\[\curl V = \partial_1 \Phi_1 \partial_2 \Phi_2 - \partial_2 \Phi_1 \partial_1 \Phi_2 = \det \begin{pmatrix} \partial_1 \Phi_1 & \partial_2 \Phi_1 \\ \partial_1 \Phi_2 & \partial_2 \Phi_2 \end{pmatrix} = \det \Jac\Phi.\]
Note that the hypothesis $\Phi \in C^2$ is doing real work here: without equality of mixed partials
the identity is simply false.

  \item Write $V_\Phi := \Phi_1 \nabla \Phi_2$ and $V_\Psi := \Psi_1 \nabla \Psi_2$. By part
\textbf{(a)}, $\curl V_\Phi = \det \Jac\Phi$ and $\curl V_\Psi = \det \Jac\Psi$, so Green's theorem
(\cref{thm:green_theorem}) turns both integrals into boundary integrals:
\[\int_\Omega \det \Jac\Phi \, dx = \int_{\partial \Omega} V_\Phi \cdot \tau \, dL,
\qquad
\int_\Omega \det \Jac\Psi \, dx = \int_{\partial \Omega} V_\Psi \cdot \tau \, dL.\]
It remains to see that the two boundary integrands agree. Fix a point of $\partial\Omega$ and let
$\tau$ be the unit tangent there. Since $\Phi = \Psi$ on all of $\partial\Omega$, the two maps agree
along the boundary curve, and therefore so do their derivatives \emph{in the direction $\tau$}:
\[\nabla \Phi_2 \cdot \tau = \partial_\tau \Phi_2 = \partial_\tau \Psi_2 = \nabla \Psi_2 \cdot \tau .\]
(The full gradients need not agree, since nothing constrains the normal derivative; but the line
integral only ever sees the tangential part.) Together with $\Phi_1 = \Psi_1$ on $\partial\Omega$
this gives $V_\Phi \cdot \tau = V_\Psi \cdot \tau$ pointwise on $\partial\Omega$, and hence
\[\int_\Omega \det \Jac\Phi \, dx = \int_\Omega \det \Jac\Psi \, dx.\]
\end{enumerate}
\begin{remark}
The statement is a genuinely striking one: $\int_\Omega \det\Jac\Phi$ is an interior quantity, the
signed area swept out by $\Phi$, yet it depends on nothing but the boundary values of $\Phi$. That
is the same phenomenon as a conservative field having path-independent work, one dimension up.
\end{remark}
\end{exercisesolution}

\begin{exercisesolution}[Solution of \cref{ex:angle_field_true_false}]
% Generator: Claude Opus 5 (High)
Begin with the computation every part relies on. Writing $X = (P,Q)$ with $P = -y/(x^2+y^2)$ and
$Q = x/(x^2+y^2)$,
\[\partial_x Q = \frac{(x^2+y^2) - 2x^2}{(x^2+y^2)^2} = \frac{y^2-x^2}{(x^2+y^2)^2},
  \qquad
  \partial_y P = \frac{-(x^2+y^2) + 2y^2}{(x^2+y^2)^2} = \frac{y^2-x^2}{(x^2+y^2)^2},\]
so $\curl X = \partial_x Q - \partial_y P = 0$ throughout $U$. The field is closed. Whether it is
also a gradient is precisely what the domain decides.
\begin{enumerate}[label=\textbf{(\alph*)}]
  \item \textbf{False.} Suppose such a $u$ existed. Then the work of $X$ around any closed loop in
    $U$ would vanish. But along the unit circle $\sigma(t) := (\cos t, \sin t)$, $t \in [0,2\pi]$,
    we have $X(\sigma(t)) = (-\sin t, \cos t) = \dot\sigma(t)$, so the integrand is
    $\lvert\dot\sigma\rvert^2 = 1$ and
\[\int_\sigma X \cdot d\sigma = \int_0^{2\pi} 1 \, dt = 2\pi \neq 0 .\]
    So no potential exists on $U$. Being closed is not enough; $U$ is not simply connected, and this
    field is the standard witness that the Poincaré lemma (\cref{thm:poincare_lemma}) genuinely needs
    that hypothesis.
  \item \textbf{True.} The half-plane $V$ is convex, hence simply connected, so a closed field on it
    is exact (\cref{cor:potentials_simply_connected}). Explicitly, $v(x,y) := \arctan(y/x)$ is $C^1$
    on $V$ --- note $x > 0$ there, so the quotient is defined and smooth --- and
\[\nabla v = \frac{1}{1 + y^2/x^2}\left(-\frac{y}{x^2},\; \frac{1}{x}\right)
   = \frac{1}{x^2+y^2}\,(-y,\, x) = X .\]
    The potential is the polar angle, which is single-valued on the half-plane but not on $U$: that
    is \textbf{(a)} and \textbf{(b)} in one sentence.
  \item \textbf{False.} Injectivity on $[0,1)$ makes $\gamma$ a simple closed curve, but it says
    nothing about \emph{where} that curve runs. Two things can go wrong. First, the loop need not
    enclose the origin at all: the circle of radius $\tfrac12$ about $(0,1)$ is a simple closed curve
    in $U$ through $(0,1)$, and since it bounds a disc on which $X$ is defined and closed, Green's
    theorem (\cref{thm:green_theorem}) gives $\int_\gamma X \cdot d\gamma = 0$. Second, even a loop
    that does enclose the origin may be traversed clockwise, giving $-2\pi$. The correct general
    statement is that the integral equals $2\pi$ times the winding number of $\gamma$ about the
    origin, and \qt{simple closed curve} pins that number down only up to sign, and only once the
    origin is known to be inside.
\end{enumerate}
\begin{remark}
This one field carries most of the chapter's content. It is closed but not exact, so it separates the
two notions; the failure is measured by an integer, the winding number; and the obstruction is a
property of the \emph{domain}, not of the field --- restrict it to any simply connected piece of $U$,
as in \textbf{(b)}, and a potential appears at once.
\end{remark}
\end{exercisesolution}
